(1) 定义 $g=(1+iC)/\tau_0^2$，原始频谱 $\widetilde A(0,\Omega)\propto e^{-\Omega^2/(2g)}$。经过色散因子后变成 $e^{-\Omega^2(1/g-i\beta_2L)/2}$；逆傅里叶得 $A(L,t)=A_0(1-i\beta_2Lg)^{-1/2}\exp[-g_Lt^2/2]$，其中 $\boxed{g_L=g/(1-i\beta_2Lg)}$，前因子保证能量归一。

(2) 令 $d=\beta_2L/\tau_0^2$，可算得 $\operatorname{Re}g_L=1/[\tau_0^2((1+Cd)^2+d^2)]$。因此 $\boxed{\tau_L=\tau_0\sqrt{(1+Cd)^2+d^2}}$，输出峰功率与输入峰功率之比为 $\tau_0/\tau_L$。频域相位模长为1，Parseval定理保证 $\int|A(L,t)|^2dt$ 完全不变。

(3) 最小化 $f(d)=1+2Cd+(1+C^2)d^2$，得到 $\boxed{d_*=-C/(1+C^2)}$、$\boxed{\tau_{\min}=\tau_0/\sqrt{1+C^2}}$。因此最佳介质积累色散量为 $\beta_2L=-C\tau_0^2/(1+C^2)$，可使峰功率提高 $\sqrt{1+C^2}$ 倍，但需要允许选择所需符号的色散。

(4) $C=0$ 时 $\tau_L=\tau_0\sqrt{1+d^2}$，任何非零纯二阶色散都使原本变换极限的脉冲变宽。若介质只能提供与最佳 $d_*$ 相反的色散符号，则无法压缩。频谱光强 $|\widetilde A|^2$ 不因乘一个单位模二次相位而改变，时间压缩来自频率分量的相位重排。
